%%%%%%%%%%%%%%%%%%%%%%%%%%%%%%%%%%%%%%%%%%%%%%%%%%%%%%%%%%%%%%%%%%%%%%%%%%%%%%%%
%% 
%% Abstract of your thesis in default document language
%% 

\cleardoubleoddpage%  Make sure to start abstract on a new odd page
\begin{abstract*}
This report documents the practical part of a seminar in AI centered on
\emph{Optimal Deep Learning of Holomorphic Operators Between Banach
Spaces} by Adcock, Dexter and Moraga~\cite{adcock2024holomorphic}
(arXiv:2406.13928). We reproduce two of the paper's three numerical
experiments -- the parametric elliptic diffusion equation (a
Hilbert-valued operator) and the Navier--Stokes--Brinkman equations (a
Banach-valued operator) -- end to end, rather than with a simplified
proxy: mixed finite element data generation in FEniCS matching the
paper's variational formulations, deterministic Clenshaw--Curtis
sparse-grid test-error evaluation, and the paper's stated training
protocol (architecture family, initialization, optimizer, and
checkpoint/restore rule). We implement and train both experiments in
PyTorch and JAX on identical data, giving a controlled comparison of the
two frameworks in addition to the comparison against the original
paper's reported results. Every point at which our implementation
necessarily deviates from an underspecified or ambiguous detail in the
paper -- for example the exact finite element family used for the
Navier--Stokes--Brinkman system, which is not stated in the paper itself
-- is documented explicitly.
\end{abstract*}

%% 
%%%%%%%%%%%%%%%%%%%%%%%%%%%%%%%%%%%%%%%%%%%%%%%%%%%%%%%%%%%%%%%%%%%%%%%%%%%%%%%%


%%%%%%%%%%%%%%%%%%%%%%%%%%%%%%%%%%%%%%%%%%%%%%%%%%%%%%%%%%%%%%%%%%%%%%%%%%%%%%%%
%% 
%% Abstract of your thesis in an additional language
%% 

\cleardoubleoddpage%  Make sure to start abstract on a new odd page
\begin{otherlanguage}{ngerman}
\begin{abstract*}
Dieser Bericht dokumentiert den praktischen Teil eines KI-Seminars zum
Thema \emph{Optimal Deep Learning of Holomorphic Operators Between Banach
Spaces} von Adcock, Dexter und Moraga~\cite{adcock2024holomorphic}
(arXiv:2406.13928). Wir reproduzieren zwei der drei numerischen
Experimente des Papers -- die parametrische elliptische
Diffusionsgleichung (ein Hilbert-wertiger Operator) und die
Navier--Stokes--Brinkman-Gleichungen (ein Banach-wertiger Operator) --
vollständig und nicht anhand einer vereinfachten Approximation: Die
Trainingsdaten werden mittels gemischter Finite-Elemente-Methoden in
FEniCS gemäß den Variationsformulierungen des Papers erzeugt, der
Testfehler wird über eine deterministische Clenshaw--Curtis-Dünngitter-Quadratur
ausgewertet, und das im Paper beschriebene Trainingsprotokoll
(Architekturfamilie, Initialisierung, Optimierer und
Checkpoint/Restore-Regel) wird eingehalten. Beide Experimente werden
sowohl in PyTorch als auch in JAX auf identischen Daten trainiert, was
einen kontrollierten Vergleich der beiden Frameworks zusätzlich zum
Vergleich mit den ursprünglich berichteten Ergebnissen ermöglicht. Jede
Stelle, an der unsere Implementierung notwendigerweise von einem im
Paper nicht eindeutig spezifizierten Detail abweicht -- etwa die genaue
Finite-Elemente-Familie für das Navier--Stokes--Brinkman-System, die im
Originalpaper nicht angegeben ist -- wird explizit dokumentiert.
\end{abstract*}
\end{otherlanguage}

%% 
%%%%%%%%%%%%%%%%%%%%%%%%%%%%%%%%%%%%%%%%%%%%%%%%%%%%%%%%%%%%%%%%%%%%%%%%%%%%%%%%
